\section{Discussion}\label{sec:discussion}

\paragraph{Evaluation protocol matters as much as architecture.}
The legacy audit shows that the same MobileViT-S drops from 88.16\% to 83.68\% Macro-F1 when
duplicate groups are kept within a single fold, although part of that gap is also due to a shorter
training schedule. Results on this dataset that use random image-level splits, including our own
legacy numbers and single-split results in the literature~\cite{hoang2026asmoe,saleh2025hybrid},
are therefore not directly comparable with the group-aware estimates in
Table~\ref{tab:main_results}. We recommend that future work on this collection release its fold
assignment or adopt the grouped split used here, which is published with the code.

\paragraph{Where the mixture helps.}
The largest F1 gains over the legacy model are on Pest (\DeltaFOnePest\,pp), Nematode
(\DeltaFOneNematode\,pp) and Healthy (\DeltaFOneHealthy\,pp), and the Pest/Fungi confusion,
the main error source of the legacy model, falls by about two thirds (145 to 46 pooled errors).
Bacteria was already near ceiling, and Phytophthora barely improves (\DeltaFOnePhytopthora\,pp)
because its confusion with Fungi persists. The routing weights (Fig.~\ref{fig:router}) are
strongly class-dependent but do not follow the intended texture-versus-global split, so the gains
cannot yet be attributed to expert specialisation. Part of the improvement also comes from data
cleaning, higher resolution and the new training recipe, which changed together with the
architecture.

\paragraph{Limitations.}
\begin{itemize}
  \item The reported gain is that of the complete pipeline; quantifying the contribution of each
  component is left for future work.
  \item Model selection uses the validation fold (early stopping), so the cross-validated estimate
  is selection-informed. An external test set from another region or season would give an unbiased
  estimate.
  \item Each validation fold contains only 11--17 Nematode images, so Nematode F1 has a standard
  error of several percentage points. Differences below about 5\,pp on this class should not be
  over-interpreted.
  \item The ROI mask is a colour heuristic. Promptable segmentation models may give cleaner crops on
  cluttered backgrounds.
  \item ASMOE~V2 has about 160\,M parameters, against 4.9--10.5\,M for the legacy models. Knowledge
  distillation into a compact student is needed for on-device use.
\end{itemize}
